\documentclass[letter,11pt]{article}
\usepackage[spanish, mexico, activeacute]{babel}
\usepackage{amsmath,amsfonts,amssymb}
\usepackage{fancyhdr}
\usepackage{fancyvrb}
\usepackage{xy}
\usepackage{graphicx}
\usepackage{latexsym}
\usepackage{enumerate}
\usepackage{tabularx}
\usepackage{multirow}
\usepackage{multicol}
\usepackage{hyperref}

\setlength{\oddsidemargin}{-0.5cm}
\setlength{\evensidemargin}{0cm}
\setlength{\textwidth}{17.5cm}
\setlength{\textheight}{31cm}
\setlength{\topmargin}{-2.5cm}

\font\bff=cmbx10 at 10truept
\font\lg=cmdunh10 at 10truept
\font\bl=cmss10 at 10truept

\newcommand{\matlab}{{\sc Matlab} }

\newcommand{\header}{\noindent%
\begin{minipage}{0.09\textwidth}
\includegraphics[width=\textwidth]{logoudec.eps}
\end{minipage}
\hspace{3mm}
\begin{minipage}{0.5\textwidth}
{\lg UNIVERSIDAD DE CONCEPCI\'ON}\hfill
\vskip-4truept
\noindent{\bff FACULTAD DE CIENCIAS}\hfill
\vskip-4truept
\noindent{\bff F\'ISICAS Y MATEM\'ATICAS}\hfill
\vskip-4truept
\noindent{\bl DEPARTAMENTO DE INGENIER\'IA MATEM\'ATICA}\hfill
\vskip4truept\hrule\hrule\vskip4truept
\par
\end{minipage}
}

\newcommand{\respuesta}[1]{
%\noindent\makebox[\textwidth][r]{
\fbox{
\begin{minipage}{\linewidth}
\hfill
\vspace{#1}
\end{minipage}
}
%}
}

\newcommand{\puntaje}[1]{\textbf{(#1 \%)}}

\newcommand{\titulo}[2]{
\newpage
\header
\begin{minipage}{0.4\textwidth}
	\begin{center}
    \ifnum#2=2
       \begin{tabular}{||p{2cm}|p{2.1cm}||}
       \hline
       Pregunta a &  Pregunta b   \\
       \hline
       \vspace{1cm}\null &     \\
       \hline
       \end{tabular}
    \else
    	\ifnum#2=3
       \begin{tabular}{||p{1.8cm}|p{1.8cm}|p{1.8cm}||}
       \hline
       Preg. a &  Preg. b &  Preg. c   \\
       \hline
       \vspace{1cm}\null &    & \\
       \hline
       \end{tabular}
     \else
       \ifnum#2=4
       \begin{tabular}{||p{1cm}|p{1cm}|p{1cm}|p{1cm}||}
       \hline
       Preg. a &  Preg. b &  Preg. c & Preg. d  \\
       \hline
       \vspace{1cm}\null &    & & \\
       \hline
       \end{tabular}
    \else
       \ifnum#2=5
       \begin{tabular}{||p{0.8cm}|p{0.8cm}|p{0.8cm}|p{0.8cm}|p{0.8cm}||}
       \hline
       Preg. a &  Preg. b &  Preg. c & Preg. d &Preg. e \\
       \hline
       \vspace{1cm}\null &    & & &\\
       \hline
       \end{tabular}
     \else
       \begin{tabular}{||p{2cm}||}
       \hline
       Pregunta a   \\
       \hline
       \vspace{1cm}\null      \\
       \hline
       \end{tabular}
    \fi
    \fi
    \fi
    \fi
  \end{center}
\end{minipage}
\begin{center}
\textbf{\Large TEST 2 VERSI\'ON #1 C\'ALCULO NUM\'ERICO 521230}\\
\vspace{5mm}
 \begin{tabular}{p{0.6\textwidth}p{0.35\textwidth}}
	\textbf{Nombre:}    &\textbf{Carrera:}\\ \\
	\textbf{Ayudante:} & \textbf{Matr\'icula:}
 \end{tabular}
 \\
\end{center}
 
 Enviar los programas solicitados en el formato solicitado al correo informado por el ayudante de su secci\'on y con copia a \textbf{numerico@ing-mat.udec.cl}.
 }
 
\usepackage[numbered,framed]{matlab-prettifier}
\lstMakeShortInline"
\lstset{
  style              = Matlab-editor,
  escapechar         = ",
  mlshowsectionrules = true,
}

 
 \pagestyle{empty}


\begin{document}
\titulo{1}{2} 				
\begin{enumerate}

%FMP
\item Considere los conjuntos de polinomios
$$
P_1=\{x^2+2x+1,x+1,x\}, \quad
P_2=\{ax^3+bx^2+x: a+b=0\}.
$$
\begin{enumerate}
  \item \puntaje{10} Construya un programa de \octave que grafique en un mismo gr\'afico todos los polinomios de $P_1$.

 \item \puntaje{20} Construya un programa de \octave que decida si existe o no un polinomio interpolante en $P_2$ de los pares ordenados
 $$
 \{(-2,-38),(-1,-7),(0,0),(1,1),(2,14)\}
 $$
\end{enumerate}

% FV
\item En dos ruteros de \octave resuelva los siguientes problemas.
\begin{enumerate}
    \item \puntaje{15} Determine el valor de $x>0$ tal que 
    \begin{equation*}
        \int_{0}^{x} t e^{t^2}\text{d}t = x.
    \end{equation*}
    \item \puntaje{15} Considere la funci\'on $f(x) = -\dfrac{1}{2}\sqrt{5x} + \sin(3x) + 1, \ x\in [1,4]$. Hallar los puntos en que el gr\'afico de $f$ tiene m\'aximos y m\'inimos locales.
\end{enumerate}
\end{enumerate}

\textbf{Desarrollo:}
\begin{enumerate}
    \item El programa solicitado en la parte a) debe tener instrucciones similares a
    
    \item
\end{enumerate}

\end{enumerate}
\titulo{2}{2}
\begin{enumerate}
%FMP
\item Descargue el archivo disponible en
\begin{center}
   \url{http://goo.gl/cQetKF}
\end{center}
este contiene una celda en donde se registran algunos precios diarios de cierre y la cantidad de acciones tranzadas de la sociedad an\'onima Aesgener en la bolsa de Santiago.

Usando estos datos y en un rutero de \octave
\begin{enumerate}
   \item \puntaje{5} Cargue los datos usando la funci\'on \texttt{load}. ?`C\'omo se llama la variable que tiene los datos?, ?`Cuantas filas y columnas tiene esta variable?.
   \item \puntaje{10} Calcule en cuantos d\'ias el precio de cierre de la acci\'on de Aesgener ha estado sobre el promedio.
   \item \puntaje{15} Grafique usando una l\'inea roja con puntos verdes el precio de cierre para cada d\'ia.
\end{enumerate}

%MS
\item 
Se quiere interpolar la funci\'on $f(x) = 1/(x^2+1)$, $x\in [-5,5]$. Para ello se consideran las dos siguientes alternativas.
\begin{enumerate}
    \item \puntaje{10}  Considere el conjunto de puntos $x_i = -5+0.5(i-1)$ con $i=1, ...,21$.
    Construya un programa de \octave que calcule el prolinomio interpolante de la funci\'on $f$ en estos puntos. Adem\'as, en un mismo gr\'afico, grafique la funci\'on $f$, los puntos de interpolaci\'on y el polinomio de interpolaci\'on obtenido.
    \item  \puntaje{15} Considere ahora el conjunto de puntos $z_i =  5\cos\left(\dfrac{(2i-1)\pi}{42}\right)$ con $i=1, ...,21$.
    Construya un programa de \octave que interpole la funci\'on $f$ en estos puntos. Adem\'as, en un mismo gr\'afico, grafique la funci\'on $f$, los puntos de interpolaci\'on y el polinomio de interpolaci\'on obtenido. Observaci\'on: los putnos $z_i$ se conocen como nodos de Chebyshev.
    \item \puntaje{5} ?`Qu\'e fen\'omeno observa en las gr\'aficas?. Justifique.
\end{enumerate}
\end{enumerate}

\textbf{Desarrollo:}
\begin{enumerate}
    \item 
    \item
\end{enumerate}

\titulo{3}{2}
\begin{enumerate}

%FMP
\item Considere el conjunto de matrices
$$
M=\left\{
\begin{bmatrix}
(-1)^i & i \\
2^{-i} & \frac{1}{i}
\end{bmatrix}
\right\}_{i=1}^{101}
$$
\begin{enumerate}
    \item \puntaje{5} Cree una funci\'on de \octave llamada \texttt{generadorM} que tenga como entrada el \'indice $i$ y como salida la i-\'esima matriz del conjunto anterior.
    \item \puntaje{10} En un rutero de \octave grafique el determinante de cada matr\'iz de $M$ versus el \'indice de estas matrices.
    \item \puntaje{15} En un rutero de \octave usando la funci\'on \texttt{det} y mediante un ciclo de iteraci\'on decida si las 101 matrices de este conjunto son o no invertibles.
\end{enumerate}

%MS
\item
Considere los siguientes datos de temperaturas en Concepci\'on  el jueves pasado.
La hora $0$ corresponde a las 0:00 hrs.
\begin{center}
\begin{tabular}{c|ccccccccccccccccccccccccccccccccc}
Hora (H) & 0 & 2 & 4 & 6 & 8 &10 &12 &14 &16 &18 &20 &22 &24 \\
\hline 
Temperatura ($^\circ C$)
& 12 & 12 & 11 &10 &10 &14 &15 & 16 &15 &14 &11& 11 & 9
\end{tabular}
\end{center}

\begin{enumerate}
    \item  \puntaje{10} Construya un programa de \octave  que ajuste los datos por el siguente modelo:
    $$
    \displaystyle T(H) =  a_0 + a_1 H + a_2H^2.
    $$
    Graficar los puntos y la curva obtenida.  Adem\'as calcular $T(9)$.
    \item  \puntaje{15}  Construya un programa de \octave  que ajuste los datos por el siguente modelo:
    $$
    \displaystyle f(H) =  c_1 + c_2 \cos\bigg(\frac{2\pi}{24} H\bigg) + c_3 \sen\bigg(\frac{2\pi}{24} H\bigg).
    $$
    Graficar los puntos y la curva obtenida.  Adem\'as calcular $f(9)$.
    \item  \puntaje{5} ?`Cu\'al de los dos modelos ajusta mejor a los datos?. Justifique.
\end{enumerate}
\end{enumerate}
\textbf{Desarrollo:}
\begin{enumerate}
    \item 
    \item
\end{enumerate}


\titulo{4}{2}
\begin{enumerate}

%MS
\item 
Considere los siguientes datos de temperaturas en Concepci\'on  el jueves pasado.
 La hora $0$ corresponde a las 0:00 hrs.
 
 \begin{center}
\begin{tabular}{c|ccccccccccccccccccccccccccccccccc}
Hora (H) & 0 & 2 & 4 & 6 & 8 &10 &12 &14 &16 &18 &20 &22 &24 \\
\hline 
Temperatura ($^\circ C$)
& 12 & 12 & 11 &10 &10 &14 &15 & 16 &15 &14 &11& 11 & 9
\end{tabular}
\end{center}
  \begin{enumerate}
\item  \puntaje{15} Construya un programa de \octave  que ajuste los datos por el siguente modelo:
$$
\displaystyle f(H) =  c_1 + c_2 \cos\bigg(\frac{2\pi}{24} H\bigg) + c_3 \sen\bigg(\frac{2\pi}{24} H\bigg).
$$
Graficar los puntos y la curva obtenida.  Adem\'as calcular $f(9)$.

\item  \puntaje{10} Ajuste los datos utilizando splines c\'ubica natural. Graficar los puntos y la curva obtenida.  Adem\'as calcular el valor del spline en  $H =9$.

\item  \puntaje{5} ?`Cu\'al de los dos modelos ajusta mejor a los datos?. Justifique.
  \end{enumerate}
% FV
\item
Considere la funci\'on definida por $f(x) = \sin(3x) + x^2, \ x\in [-\pi,\pi]$, y la recta $L$ de ecuaci\'on $2x-y+1=0$. Construya un programa de \octave que:
\begin{enumerate}
    \item \puntaje{10} Grafique, en un mismo gr\'afico (con distinto color), la funci\'on $f$ y la recta $L$.
    \item \puntaje{20} Encuentre los puntos del intervalo $[-\pi,\pi]$ en que la recta tangente al gr\'afico de $f$ es paralela a la recta $L$.
\end{enumerate}
\end{enumerate}

\textbf{Desarrollo:}
\begin{enumerate}
    \item 
    \item
\end{enumerate}
\titulo{5}{2}
\begin{enumerate}

%FMP
\item El n\'umero $6174$ es conocido como una constante de \textbf{Kaprekar}. Todos los n\'umeros naturales menores que $10000$ y mayores que $1000$, a excepcion de los que tengan sus digitos repetidos, se pueden transformar en $6174$ ordenando los d\'itigos del n\'umero en orden descendente y ascendente, luego calcular la resta entre estos dos n\'umeros y finalmente repitiendo este proceso.

Por ejemplo, para el n\'umero $n=2368$, sus digitos ordenados en orden ascendente son $2368$ y descendente son $8632$. As\'i la diferencia entre ellos es
$$
8632-2368=6264,
$$
repitiendo este proceso con el n\'umero reci\'en calculado
$$
6642-2466=4176
$$
similarmente, se llega a constante de Kaprekar.
$$
7641-1467=6174.
$$
As\'i, se dice que el n\'umero $n=2368$ tiene 3 pasos de Kaprekar.

\begin{enumerate}
    \item \puntaje{20} Cree una funci\'on de \octave que tenga como entrada un n\'umero natural menor que $10000$ y mayor que $1000$ y que tenga como salida los pasos de Kaprekar de dicho n\'umero.
    \item \puntaje{10} Grafique los pasos de Kaprekar versus el n\'umero para los n\'umeros naturales entre $6000$ y $6200$.
\end{enumerate}



% FV
\item
Considere la funci\'on definida por $f(x) = \sin(x^2) - e^{3x+1}, \ x\in [-\pi,0]$, y la recta $L$ de ecuaci\'on $2x+y+4=0$. Construya un programa de \octave que:
\begin{enumerate}
    \item \puntaje{10} Grafique, en un mismo gr\'afico (con distinto color), la funci\'on $f$ y la recta $L$.
    \item \puntaje{20} Encuentre los puntos del intervalo $[-\pi,0]$ en que la recta $L$ intersecta el gr\'afico de $f$.
\end{enumerate}
\end{enumerate}

\textbf{Desarrollo:}
\begin{enumerate}
    \item 
    \item
\end{enumerate}

\end{document}
